% ==================================================================
\label{sec:system}
% ==================================================================

\subsection{Islanded Inverter-Based Microgrid}
\label{ssec:islanded_microgrid}

We consider the islanded operation of the IEEE 123-bus distribution
feeder with a fully inverter-based distributed energy resource (DER)
portfolio. The controllable fleet comprises distributed photovoltaics
(DPVs), battery energy storage systems (BESSs), and electric-vehicle
charging stations with vehicle-to-grid capability (EVCS/V2G). No
bulk-grid slack bus is available; voltage and frequency references are
established by grid-forming (GFM) inverter resources deployed inside the
microgrid. The resulting system is a low-inertia, topology-dependent
feeder in which fast frequency response (FFR) must be delivered by
geographically distributed inverter-interfaced assets.

The feeder is partitioned into VPP zones for operational coordination.
Let $\mathcal{Z}$ denote the set of VPP zones with $|\mathcal{Z}|=3$,
$\mathcal{A}_z$ the set of controllable DER agents associated with zone
$z$, and $\mathcal{A}$ the full controllable agent set
\begin{equation}
    \mathcal{A}=\bigcup_{z\in\mathcal{Z}}\mathcal{A}_z,
    \qquad |\mathcal{A}|=41.
\end{equation}
The implemented ancillary-market (AM) controller acts at DER granularity
with $|\mathcal{A}|=41$ controllable agents. These DER agents are not
treated as isolated points: their observations are embedded through the
active 123-bus feeder graph so that electrical distance, switching
status, and background grid state can influence the fast
frequency-control decision.

\subsection{Dual-Product AM Service Focus within Multi-Service VPP Operation}
\label{ssec:multiservice_vpp}

The broader VPP framework supports both energy-market (EM) and AM
service accounting. EM represents the slower service layer for
active/reactive dispatch commitments and schedule tracking, whereas AM
represents the fast service layer for reserve commitment and FFR
activation under frequency events. In this paper, the learning
controller is formulated for the AM fast loop. EM schedules, zonal
prices, topology choices, and service commitments enter as exogenous
context generated by the upper layers, while the GraphSAGE--MAPPO policy
learns how the DER fleet should deliver fast active-power support during
frequency excursions.

A distinctive feature of the implemented AM controller is that it is
formulated as a \emph{dual-product} service problem. Each controllable
DER simultaneously commits a fast active-power reference
$P^{\mathrm{ref}}_{i,t}$ (analogous to an energy-product commitment) and
a frequency-droop gain $K^{\mathrm{droop}}_{i,t}$ (analogous to a
capacity-product reservation). The fast environment couples both
commitments to the equivalent frequency dynamics, so that the policy
must trade off transient frequency support, storage-state health, and
market value within a single decision. This is structurally different
from a classical droop baseline, which fixes $K^{\mathrm{droop}}$ and
exposes only $P^{\mathrm{ref}}$ to the controller, and from a
single-action MAPPO baseline, which exposes only $P^{\mathrm{ref}}$
while a fixed droop is bypassed. Both classical-droop and single-action
MAPPO modes are retained for ablation, but the proposed AM controller
operates in dual-product mode by default.

This separation is deliberate. The earlier dual-loop EM+AM formulation
couples slow EM dispatch and fast AM response in a single learning
problem, which dilutes the learning signal for frequency control. The
revised implementation isolates the AM objective, keeps the fast
environment contract used by the simulator, and evaluates policies
using frequency-quality KPIs (maximum absolute frequency deviation,
RoCoF, violation fraction, and nadir) together with service-accounting
diagnostics for the dual-product objective. The market interpretation
is retained through Layer~1 signals and VPP membership, but the
learning problem studied here is the AM real-time execution problem.

\subsection{Layered AM Execution Architecture}
\label{ssec:layered_architecture}

The operating framework is organized into three layers. Layer~0 is the
DSO-side grid-operation layer. It monitors feeder states and network
constraints, processes topology changes and contingencies, and outputs
the active feeder graph, grid-feasibility envelope, and zonal operating
signals, denoted compactly by $G_t$, $\Omega_t^{\mathrm{grid}}$, and
$\lambda_t^{\mathrm{zone}}$. In the implementation, $G_t$ is represented
by the energized edge index of the pandapower feeder model and may
change when tie-switch configurations or contingency events alter the
active network.

Layer~1 is the VPP service-coordination layer. It converts Layer~0
signals into zone-level service context, including active/reactive
schedules, reserve obligations, VPP membership, zonal locational
marginal prices, and capacity-price coefficients. These quantities are
not optimized by the AM trainer itself; they are provided as contextual
service signals and enter the AM objective through service-accounting
terms (capacity payment proportional to droop-gain reservation, energy
revenue proportional to delivered FFR power).

Layer~2 is the real-time execution layer. It uses a feeder-level
GraphSAGE--MAPPO policy operating in dual-product AM mode. The policy
receives the active feeder graph and a full-feeder node-feature matrix,
encodes all buses through GraphSAGE, extracts the embeddings at the 41
controllable DER buses, and outputs, for each DER, a two-dimensional
normalized action $(a^{P}_{i,t},a^{K}_{i,t})$ that determines its fast
active-power reference and its local droop gain. The simulator consumes
an action contract of dimension
$2|\mathcal{A}|+|\mathcal{Z}|=2\cdot 41+3=85$. The first $|\mathcal{A}|$
entries carry the per-DER power references, the next $|\mathcal{A}|$
entries carry the per-DER droop gains, and the final $|\mathcal{Z}|$
entries are legacy VPP-aggregated slots retained for interface
compatibility (set to zero in dual-product mode, since per-DER droop
replaces VPP-mean droop aggregation). In this way, Layer~0 defines what
is grid-feasible, Layer~1 defines the service context, and Layer~2
determines how individual DERs deliver fast ancillary support in real
time.

\begin{figure}[!t]
  \centering
  \includegraphics[width=\columnwidth]
    {img/System architecture.png}
  \caption{Layered AM execution architecture for topology-aware VPP
ancillary-service delivery in an islanded inverter-based microgrid. The
GraphSAGE--MAPPO controller operates on the active feeder graph,
extracts DER-bus embeddings, and produces a per-DER dual-product action
$(a^{P}_{i,t},a^{K}_{i,t})$ that maps to a fast active-power reference
and a local droop gain.}
  \label{fig:system_arch}
\end{figure}

\subsection{Topology-Dependent FFR Deliverability}
\label{ssec:topology_deliverability}

FFR deliverability is topology-dependent in islanded feeders. A DER may
have sufficient local capability, such as BESS state of charge or
inverter headroom, but branch congestion, switching status, voltage
sensitivity, and electrical distance can limit whether its support can
be delivered without violating network constraints. Likewise, the
effectiveness of active-power support during a frequency event depends
on which buses remain electrically connected to the grid-forming
reference and how DER support propagates through the energized feeder
paths.

To represent this dependence, the time-varying feeder state is expressed
as
\begin{equation}
    G_t=(\mathcal{N},\mathcal{E}_t,\mathbf{X}_t,\mathbf{E}_t),
\end{equation}
where $\mathcal{N}$ is the bus set, $\mathcal{E}_t$ is the energized
branch set, $\mathbf{X}_t$ denotes node features, and $\mathbf{E}_t$
denotes physical edge attributes. The implemented policy uses
$\mathcal{E}_t$ through an edge-index representation and uses a
full-feeder node-feature matrix derived from frequency states, DER
states, loads, voltages, and service context. Edge attributes such as
resistance, reactance, and thermal rating remain part of the environment
and Layer~0 feasibility logic, while the GraphSAGE message function
operates on the unweighted active topology.

Topology awareness is therefore not introduced as a cosmetic graph
representation. It is required because the physical effect of the same
DER action depends on the energized feeder graph, and because held-out
switching configurations can change which neighboring buses provide
useful context for a DER's response.

\subsection{Multi-Timescale Information Flow}
\label{ssec:multitimescale_flow}

Information exchange follows distinct operational timescales. Layer~0
updates topology and grid-feasibility signals at planning or
event-triggered timescales. Layer~1 translates these updates into VPP
service context at scheduling intervals. Layer~2 operates in the fast
AM channel: it observes frequency deviation, RoCoF, DER availability
indicators (state of charge, droop-headroom utilization, recent
commitments, V2G availability proxy), VPP-zone membership, and the
active feeder graph, then produces the dual-product action for the
controllable DER fleet.

Feedback closes the loop after each disturbance response. The
environment records delivered FFR energy, frequency nadir, maximum
absolute frequency deviation, RoCoF, violation fraction, topology
changes, and feasibility diagnostics. These telemetry streams, together
with service-accounting quantities (LMP-weighted energy revenue and
droop-gain-weighted capacity revenue), are used for policy training,
hyperparameter screening, and subsequent evaluation.

\paragraph{Cooperative Primary--Secondary--FFR Coordination}
The three control layers are intentionally arranged in a
\emph{cooperative}, not a substitutive, configuration. The VPP-dispatched
FFR (Layer~2) saturates rapidly within $\sim$0.1--2~s (T$_{\text{BESS}}$=100~ms,
T$_{\text{V2G}}$=300~ms, T$_{\text{PV-curtail}}$=1~s), the GFM primary droop
(Layer~0/1 embedded in the swing dynamics) ramps in over $T_g$=1~s and
provides sustained support up to its rating, and the GFM secondary AGC
activates after a 10~s delay (event-armed PI) to drive the residual
$\Delta f$ back to nominal. Storage-based FFR is the first responder
because its energy reservoir is small but its power electronics latency
is the lowest in the system, in line with the Nordic
FCR$_{\text{I}}$ priority order for converter-interfaced reserves in
isolated grids~\cite{mota2024fcri,papakonstantinou2021bessagc}. To
prevent the AGC integrator from accumulating frequency error that the
faster FFR loop is already correcting---and the resulting AGC overshoot
once FFR fades---we freeze the secondary integrator while
$\mathbf{1}_{\text{FFR active}}=1$, following the anti-windup
strategy of~\cite{subotic2023dualportgfm}. This decoupling yields a
clean three-stage response (FFR peak $\rightarrow$ primary ramp
$\rightarrow$ AGC restoration) that we visualise empirically in
Subsection~\ref{ssec:cooperative_dispatch}.

The following
section formalizes the physical constraints and service models used by
this architecture. The AM Dec-POMDP and the feeder-level
GraphSAGE--MAPPO policy are then developed in
Sections~\ref{sec:pomdp} and~\ref{sec:method}, respectively.
